% =============================================================================================
% §6 Results — STRUCTURED STUBS ONLY. The pre-registered run has not been executed.
%
% CONTEXT-STRIPPING RULE (CLAUDE.md standing rules). Every slot below is \TBD{} or \PH{}, in red,
% and NO slot contains a plausible value. A reader who crops a paragraph out of this section must
% not be able to mistake any of it for a measurement. When the run lands, each \TBD is replaced
% from a single named field of the verdict manifest -- the field is named beside every slot so the
% substitution is mechanical and cannot be improvised.
%
% BOTH interpretation paragraphs in §6.6 are written NOW and dated, before any result exists, so
% that neither can be tuned to an outcome. That is the point of writing them here.
% =============================================================================================
\section{Results}
\label{sec:results}

\emph{The pre-registered run has not been executed. Every quantity in this section is an empty
slot bound to a named field of the verdict manifest, and every figure is drawn on placeholder data
that is implausible by construction. Nothing here is a measurement.}

\subsection{The pre-second-stage gate}
\label{sec:res-gate}

The legibility gate of \cref{sec:mech-gate} runs on real data before any second-stage training is
authorized, and it is a hard stop rather than a diagnostic: it requires at least $0.90$ sign
accuracy on \emph{every one} of the six live microstructure dimensions, with no mean clause and no
floor. Its outcome determines whether the rest of this section exists.

\begin{center}
\begin{tabular}{ll}
\toprule
quantity & value \\
\midrule
per-dimension sign accuracy, dims 7--12 & \TBD{six values} \\
minimum across the six & \TBD{min} \\
gate outcome & \TBD{pass or hard stop} \\
\bottomrule
\end{tabular}
\end{center}
% manifest binding: gates.gate2_legibility.{new_by_seed, min, threshold=0.90, pass}

\subsection{Primary outcome}
\label{sec:res-primary}

The five clauses of \cref{sec:decision} are evaluated conjunctively: the primary outcome is
\textsc{survives} only if none fails, and \textsc{null} otherwise. Two guards sit outside the
clauses and are \textsc{halt}-only. \Cref{fig:verdict} draws the surface.

\begin{center}
\begin{tabular}{lll}
\toprule
clause & tests & outcome \\
\midrule
1 paired CI & $\Delta\IR(4-5)$ lower bound $> 0$ & \TBD{} \\
2 MDE & $\Delta\IR(4-5) \ge$ paired MDE & \TBD{} \\
3 placebo validity & $\IR(4) - \IR(2)$ lower bound $> 0$ & \TBD{} \\
4 economic floor & materiality against the $0.5$ floor & \TBD{} \\
5 deflated Sharpe & against $SR_0$ at $N = 60$ trials & \TBD{} \\
\midrule
degeneracy guard & constant-direction book, any cell & \TBD{halt or clear} \\
power guard & across-seed range vs claimed effect & \TBD{halt or clear} \\
\midrule
\textbf{primary} & & \TBD{SURVIVES / NULL / INCONCLUSIVE} \\
\bottomrule
\end{tabular}
\end{center}
% manifest binding: clauses.{1_paired_ci, 2_mde_paired, 3_placebo_validity, 4_economic_floor,
% 5_dsr}.passed; degeneracy_guard.halted; power_guard.halted; primary. The rule STRINGS printed in
% Fig 8 are read from the manifest, not restated.

Both specifications required by the amendment window are reported from the same artifacts, and
their agreement or disagreement is stated here rather than in a footnote.

\begin{center}
\begin{tabular}{lll}
\toprule
 & v1.5 specification (primary) & retained alternative \\
\midrule
primary outcome & \TBD{} & \TBD{} \\
specifications agree & \multicolumn{2}{l}{\TBD{yes or no --- a first-class finding either way}} \\
\bottomrule
\end{tabular}
\end{center}
% manifest binding: dual_specification.* — a manifest lacking this field cannot be quoted (§7 v1.5)

\subsection{Effect sizes and secondary diagnostics}
\label{sec:res-effects}

Information coefficient, rank information coefficient, mean absolute error and $R^2$ are reported
as secondary diagnostics and do not enter the decision rule.

\begin{center}
\begin{tabular}{lccccc}
\toprule
cell & 1 BSQ/OHLCV & 2 FSQ/OHLCV & 3 BSQ/+micro & 4 FSQ/+micro & 5 placebo \\
\midrule
net IR, seed mean & \TBD{} & \TBD{} & \TBD{} & \TBD{} & \TBD{} \\
across-seed range & \TBD{} & \TBD{} & \TBD{} & \TBD{} & \TBD{} \\
RankIC & \TBD{} & \TBD{} & \TBD{} & \TBD{} & \TBD{} \\
decision activity at $\kappa^\ast$ & \TBD{} & \TBD{} & \TBD{} & \TBD{} & \TBD{} \\
\bottomrule
\end{tabular}
\end{center}
% manifest binding: per_cell.{ir_seed_mean, ir_range_across_seeds, rankic}; degeneracy_guard.
% per_cell.activity_decisions_seed_mean. The activity row is printed whatever it reads, per §4.4.

The three marginals involving a BSQ arm are reported and not claimed, for the reason
\cref{sec:lim-baseline} gives, and the required disclosure ships with them verbatim.

\subsection{Cost stress and seed dispersion}
\label{sec:res-stress}

\Cref{fig:cost} sweeps the round-trip cost around the pinned $0.30\%$; \cref{fig:seeds} plots each
cell's five seeds against the size of the effect being claimed. Both figures are built against the
manifest schema and are watermarked until the run lands.

\subsection{Interpretation}
\label{sec:res-interpretation}

\emph{Both paragraphs below were written on 2026-08-10, before the run was executed and before any
result existed. Exactly one will be retained; the other is deleted rather than softened. They are
placed here, in advance, so that neither can be written to fit an outcome.}

\paragraph{If the outcome is \textsc{survives}.} A positive $\Delta\IR(4-5)$ that clears the
paired confidence interval, the minimum detectable effect, the economic floor and the deflation
benchmark establishes that real trade-flow imbalance carries information beyond input capacity, and
that a tokenizer built to preserve it delivers that information to a cost-aware trading decision.
The claim ends there, and three limits bind it immediately. It is not a calibrated estimate of how
much: \cref{sec:lim-placebo} states that $\Delta\IR(4-5)$ mixes information with an unquantified
capacity handicap, so the measured difference is an upper bound on the information effect rather
than an estimate of it. It is not evidence that our BSQ baseline is competitive, so the
quantizer-axis marginals remain reported rather than claimed. And it is one exchange, one venue,
one instrument class and one year of scored data, with no regime-transfer test --- a positive
result here is an existence proof on this slice, not a property of markets. The mechanism of
\cref{sec:mechanism} would remain a fixture result; a positive outcome is consistent with it and
does not confirm it, because many differences between cells 2 and 4 could produce the same sign.

\paragraph{If the outcome is \textsc{null} or \textsc{inconclusive}.} This is the outcome the
design's own arithmetic makes most likely, and we pre-committed it as publishable before any result
existed, so it is reported without reframing. A null here means that under this pre-registered
decision rule, at this budget, on this slice, the microstructure arm did not beat its placebo by a
margin the design could resolve. It does \emph{not} mean microstructure carries no information: the
minimum detectable effect is $3.518$ annualized IR against a prior of $|\text{RankIC}| = 0.027$, so
the design would fail to detect a real effect of the size we actually expect, and
\cref{sec:lim-power} states that the effect we can detect is larger than the effect we predict.
Nor does it mean the mechanism of \cref{sec:mechanism} is wrong --- the mechanism is established on
a fixture and is not contingent on this outcome. What a null does establish is narrower and still
worth publishing: that the pipeline, with the interface defect removed and the legibility gate
passed, does not convert this microstructure signal into net information ratio at this scale.
Given how much of this literature reports only positive results, a pre-registered null with its
power analysis stated in advance is a contribution to the record rather than a failed experiment.
\Cref{sec:lim-power} names what would change the answer: more scored data, a lower-cost venue, or
an effect larger than our own prior.
